\section{PDO: Koneksi dan Prepared Statement}

\textbf{PDO} (PHP Data Objects) adalah abstraksi database di PHP yang mendukung \textbf{banyak DBMS} (MySQL, PostgreSQL, SQLite, SQL Server, dll.) melalui interface yang seragam. Berbeda dengan MySQLi yang hanya untuk MySQL, PDO memungkinkan Anda berpindah database dengan perubahan konfigurasi minimal. PDO hanya mendukung gaya OOP dan selalu menggunakan prepared statement secara alami. Banyak framework PHP modern (Laravel, Symfony) menggunakan PDO sebagai lapisan database \cite{phpManual}.

\begin{table}[ht]
\centering
\begin{tabular}{|P{3.5cm}|P{5cm}|P{5cm}|}
\hline
\textbf{Aspek} & \textbf{MySQLi} & \textbf{PDO} \\
\hline
Database & Hanya MySQL & 12+ DBMS \\
\hline
Gaya API & Prosedural + OOP & Hanya OOP \\
\hline
Named params & Tidak & Ya (\code{:nama}) \\
\hline
Prepared stmt & Ya & Ya \\
\hline
Performa MySQL & Sedikit lebih cepat & Sedikit lebih lambat \\
\hline
Portabilitas & Rendah & Tinggi \\
\hline
\end{tabular}
\caption{Perbandingan MySQLi dan PDO}
\end{table}

\begin{webcode}[language=PHP, caption={Koneksi dan CRUD dengan PDO}]
<?php
// Koneksi PDO
try {
  $pdo = new PDO(
    "mysql:host=localhost;dbname=toko_online;charset=utf8mb4",
    "root", "",
    [
      PDO::ATTR_ERRMODE            => PDO::ERRMODE_EXCEPTION,
      PDO::ATTR_DEFAULT_FETCH_MODE => PDO::FETCH_ASSOC,
      PDO::ATTR_EMULATE_PREPARES   => false
    ]
  );
} catch (PDOException $e) {
  die("Koneksi gagal: " . $e->getMessage());
}

// INSERT dengan named parameters
$stmt = $pdo->prepare("INSERT INTO produk (nama, harga) VALUES (:nama, :harga)");
$stmt->execute([':nama' => 'Mouse', ':harga' => 150000]);
echo "ID baru: " . $pdo->lastInsertId();

// SELECT
$stmt = $pdo->prepare("SELECT * FROM produk WHERE harga > :min");
$stmt->execute([':min' => 100000]);
$produk = $stmt->fetchAll(); // FETCH_ASSOC otomatis

foreach ($produk as $p) {
  echo "{$p['nama']}: Rp {$p['harga']}\n";
}
?>
\end{webcode}

